%%%==============EX_1============%%%
\begin{ex}%[2H5V1-5] 
	Trong không gian với hệ trục tọa độ $Oxyz$, cho hai điểm $A(1;2;3)$, $B(5;-4;-1)$ và mặt phẳng $(P)$ qua $Ox$ sao cho $\mathrm{d}(B,(P))=2\mathrm{d}(A,(P))$ cắt $AB$ tại $I(a;b;c)$ nằm giữa $AB$. Tính $a+b+c$.
	\choice
	{$8$}
	{$6$}
	{$12$}
	{\True $4$}
	\loigiai{
		Mặt phẳng $(P)$ qua $Ox$ nên phương trình mặt phẳng $(P)$ có dạng $by+c=0$ $(b^2+c^2>0)$.\\
		Ta có
		\begin{eqnarray*}
			&&\mathrm{d}(B,(P)) =2\mathrm{d}(A,(P))\\
			&\Leftrightarrow&\dfrac{|-4-c|}{\sqrt{b^2+c^2}} =2\cdot \dfrac{2b+3c}{\sqrt{b^2+c^2}}\\
			& \Leftrightarrow &\hoac{&-4b-4b=4b+6c\\&-4b-4b=-4b-6c}\\
			&\Leftrightarrow &\hoac{&8b+7c=0\\ &c=0.}
		\end{eqnarray*}
		\begin{itemize}
			\item Trường hợp 1: $8b+7c=0$.\\
			Chọn $b=7$, $c=-8$, khi đó $(P)\colon7y-8z=0$.\\
			Xét $f(x;y)=7y-8z$, thay tọa độ $A$, $B$ vào ta được $(7\cdot 2-8\cdot 3)(7\cdot (-4)-8\cdot (-1))>0$.\\
			Suy ra $A$, $B$ nằm cùng phía so với $(P)$ (loại).
			\item Trường hợp 2: $c=0$, suy ra phương trình $(P):y=0$.\\
			Thay tọa độ $A$, $B$ vào ta được $2 \cdot (-4) <0$.\\
			Suy ra $A$, $B$ nằm khác phía so với $(P)$. Do đó đường thẳng $AB$ cắt $(P)$ tại $I$ nằm giữa $AB$.\\
			PTTS của đường thẳng $AB\colon \heva{&x=1+4t\\ &y=2-6t\\ &z=3-4t} (t\in \mathbb{R})$.
		\end{itemize}
		Tọa độ điểm $I$ là nghiệm hệ phương trình
		$$\heva{&x=1+4t\\ &y=2-6t\\ &z=3-4t\\ &y=0} \Leftrightarrow \heva{&t=\dfrac{1}{3}\\ &x=\dfrac{7}{3}\\ & y=0\\ &z=\dfrac{5}{3}} \Rightarrow I\left( \dfrac{7}{3};0;\dfrac{5}{3} \right).$$
		Vậy $a+b+c=4$.
	}
\end{ex}

%%%==============EX_2============%%%
\begin{ex}%[2H5V2-5] 
	Trong KG $Oxyz$, cho $A(2;3;-1)$, $B(2;3;2)$, $C(-1;0;2)$.Tìm tọa độ điểm $M$ thuộc mặt phẳng $(Oxz)$ để $S=\left|\overrightarrow{MA}-4\overrightarrow{MC}\right|+\left|\overrightarrow{MA}+\overrightarrow{MB}+\overrightarrow{MC}\right|$ nhỏ nhất.
	\choice
	{\True $M\left(-1;0;\dfrac{7}{3} \right)$}
	{$M(0;3;0)$}
	{$M\left(1;0;\dfrac{7}{3} \right)$}
	{$M\left(-\dfrac{1}{2};0;2\right)$}
	\loigiai{
		Gọi $G$ là trọng tâm tam giác $ABC$, suy ra $G(1;2;1)$.\\
		Gọi $H(x;y;z)$ là điểm thỏa mãn $\overrightarrow{HA}-4\overrightarrow{HC}=\overrightarrow0$.
		$$\Leftrightarrow\heva{&{2-x=4(-1-x)} \\
			&{3-y=4(0-y)} \\
			&{-1-z=4(2-z)}
		} \Leftrightarrow \heva{&{x=-2} \\
			&{y=-1} \\
			&{z=3}
			&} \Rightarrow H(-2;-1;3).$$
		Nhận thấy $G$, $H$ nằm về hai phía đối với mặt phẳng $(Oxz)$ và $HG=\sqrt{22}$.\\
		Ta có
		\begin{eqnarray*}
			&S&=\left|\overrightarrow{MA}-4\overrightarrow{MC}\right|+\left|\overrightarrow{MA}+\overrightarrow{MB}+\overrightarrow{MC}\right|\\
			&&=\left|\overrightarrow{MH}+\overrightarrow{HA}-4\overrightarrow{MH}-4\overrightarrow{HC}\right|+\left|\overrightarrow{MG}+\overrightarrow{GA}+\overrightarrow{MG}+\overrightarrow{GB}+\overrightarrow{MG}+\overrightarrow{GC}\right|\\
			&& =\left|-3\overrightarrow{MH}\right|+\left|3\overrightarrow{MG}\right|\\
			&&=3(MH+MG)\ge 3GH=3\sqrt{22}.
		\end{eqnarray*}
		Đẳng thức xảy ra khi và chỉ khi $H$, $M$, $G$ thẳng hàng theo thứ tự.\\
		Lại do $M\in (Oxz)$ nên $S$ đạt giá trị nhỏ nhất khi và chỉ khi $M$ là giao điểm của đường thẳng $GH$ với mặt phẳng $(Oxz)$.\\
		Đường thẳng $GH$ có phương trình $\heva{&{x=1+3t} \\
			&{y=2+3t} \\
			&{z=1-2t}.}$\\
		Mặt phẳng $(Oxz)$ có phương trình $y=0$.\\
		Ta có $M\in GH\Rightarrow M(1+3t;2+3t;1-2t).$\\		
		Mặt khác $M\in (Oxz)\Rightarrow 2+3t=0\Leftrightarrow t=-\dfrac{2}{3}$.\\
		Vậy $M\left(-1;0;\dfrac{7}{3} \right)$.
	}
\end{ex}

%%%==============EX_3============%%%
\begin{ex}%[2H5C1-3] 
	Trong KG $Oxyz$ cho mặt phẳng $(P)$ đi qua điểm $M(9;1;1)$ cắt các tia $Ox$, $Oy$, $Oz$ lần lượt tại $A$, $B$, $C$ ($A$, $B$, $C$ không trùng với gốc tọa độ). Thể tích tứ diện $OABC$ đạt giá trị nhỏ nhất là bao nhiêu?
	\choice
	{\True $\dfrac{81}{2}$}
	{$\dfrac{243}{2}$}
	{$\dfrac{81}{6}$}
	{$243$}
	\loigiai{
		Giả sử $A(a;0;0)$, $B(0;b;0)$, $C(0;0;c)$ với $a$, $b$, $c$ $>0$.\\
		Mặt phẳng $(P)$ có phương trình đoạn chắn là $\dfrac{x}{a}+\dfrac{y}{b}+\dfrac{z}{c}=1$.\\
		Vì mặt phẳng $(P)$ đi qua điểm $M(9;1;1)$ nên $\dfrac{9}{a}+\dfrac{1}{b}+\dfrac{1}{c}=1$.\\
		Ta có $1=\dfrac{9}{a}+\dfrac{1}{b}+\dfrac{1}{c} \ge 3\sqrt[{3}]{\dfrac{9}{a\cdot b\cdot c}}
		\Rightarrow a\cdot b\cdot c\ge 243$.\\ 
		Khi đó $V_{OABC}=\dfrac{1}{6} a\cdot b\cdot c\ge \dfrac{243}{6}=\dfrac{81}{2}$.\\
		Vậy thể tích tứ diện $OABC$ đạt giá trị nhỏ nhất là $\dfrac{81}{2}$.
	}
\end{ex}

%%%==============EX_4============%%%
\begin{ex}%[2H5C1-5] 
	Trong KG $Oxyz$, cho các điểm $M(m; 0; 0)$, $N(0; n; 0)$, $P(0; 0; p)$ không trùng với gốc tọa độ và thỏa mãn $m^2+n^2+p^2=3$. Tìm giá trị lớn nhất của khoảng cách từ $O$ đến mặt phẳng $(MNP)$.
	\choice
	{$\dfrac{\sqrt{3}}{9}$}
	{$\dfrac{1}{3}$}
	{\True $\dfrac{\sqrt{3}}{3}$}
	{$\dfrac{1}{27}$}
	\loigiai{
		Mặt phẳng $(MNP)$ có phương trình là $\dfrac{x}{m}+\dfrac{y}{n}+\dfrac{z}{p}=1$.\\
		Theo bất đẳng thức Bunhia-Copsky ta có $$\left(m^2+n^2+p^2 \right) \left(\dfrac{1}{m^2}+\dfrac{1}{n^2}+\dfrac{1}{p^2} \right)\ge 9\Rightarrow \dfrac{1}{m^2}+\dfrac{1}{n^2}+\dfrac{1}{p^2} \ge \dfrac{9}{m^2+n^2+p^2}=3.$$
		Khi đó $\mathrm{d}\left(O;(P)\right)=\dfrac{1}{\sqrt{\dfrac{1}{m^2}+\dfrac{1}{n^2}+\dfrac{1}{p^2}}} \le \dfrac{\sqrt{3}}{3}$.\\
		Dấu bằng xảy ra khi $m=n=p=1$.\\
		Vậy khoảng cách lớn nhất từ $O$ đến $(MNP)$ bằng $ \dfrac{\sqrt{3}}{3}$.}
\end{ex}
\Closesolutionfile{ans}
\indapan{10}{ans/ans-C5B3CD5}
%%%==============EX_5============%%%
\TNSA
\Opensolutionfile{ans}[ans/ans-C5B3CD5-KQ]
\begin{ex}%[2H5V1-3]
	Trong KG $Oxyz$, cho điểm $A(2;1;-1)$. Gọi $(P)\colon  ax+by+cz=0$ là mặt phẳng chứa trục $Oy$ sao cho khoảng cách từ $A$ đến $(P)$ là lớn nhất. Tính $a+b+c$.
	\shortans{$1$}
	\loigiai{
		\immini{
			Gọi $H$ là hình chiếu vuông góc của điểm $A$ lên mặt phẳng $(P)$, $A'$ là hình chiếu vuông góc của điểm $A$ lên trục $Oy$. Suy ra $A'(0;1;0)$.\\
			Khi đó $\mathrm{d}\left(A;(P)\right)=AH \le AA'$.\\
			Độ dài đoạn thẳng $AH$ dài nhất khi $H$ và $A'$ trùng nhau. Khi đó mặt phẳng $(P)$ nhận $\overrightarrow{A'A}=(2;0;-1)$ làm véc-tơ pháp tuyến.\\
			Phương trình mặt phẳng $(P)$ đi qua $A'(0;1;0)$ có véc-tơ pháp tuyến $\overrightarrow{A'A}=(2;0;-1)$ là 
			$$2(x-0)+0(y-1)+(-1)(z-0)=0\Leftrightarrow 2x-z=0.$$
			Vậy $a+b+c=1$.	
		}{
			\begin{tikzpicture}[line join=round, line cap=round,>=stealth,thick,scale=0.5]
				%Khai báo điểm----------------------
				\path (0,0) coordinate (H);
				\path (0,3) coordinate (A);
				\path (3,-1) coordinate (A');
				\path (-1,1) coordinate (D);	
				\path (-3,-3) coordinate (E);
				\path (5,-3) coordinate (F);	
				\path (5,0) coordinate (I);
				\path (1,-2) coordinate (G);
				\path pic[draw,angle radius=5pt]{right angle= A--A'--I};
				\path pic[draw,angle radius=5pt]{right angle= A--H--A'};
				\path pic["\scriptsize$P$", angle eccentricity=0.6,draw,angle radius=20pt]{angle= F--E--D};
				%vẽ---------------------------------
				\draw (A)--(H)--(A')--cycle;
				\draw (D)--(E)--(F);
				\draw (I)--(A')--(G);
				\node at (4,0){$Oy$};
				%hiển thị tên điểm----------------------
				\foreach \t/\g in {H/-90,A'/-90,A/90}{
					\draw[fill=black] (\t) circle (1pt) node[shift={(\g:10pt)}]{$ \t $};
				}
			\end{tikzpicture}	
		}
	}
\end{ex}

%%%==============EX_6============%%%
\begin{ex}%[2H5V1-3]
	Trong KG $Oxyz$, cho điểm $A(2;1;1)$. Gọi $(P)\colon ax+by+c=0$ là mặt phẳng chứa trục $Oy$ sao cho khoảng cách từ $A$ đến $(P)$ lớn nhất. Tính $a+b+c$.
	\SA{3}
	\loigiai{
		\immini{
			Gọi $H$ và $K$ lần lượt là hình chiếu của $A$ trên $(P)$ và trục $Oy$.\\
			Ta có $\mathrm{d}\left(A,(P)\right)=AH\le AK$.\\
			Do đó khoảng cách từ $A$ đến $(P)$ lớn nhất khi $H\equiv K(0;1;0)$.\\
			Phương trình $(P)$ đi qua $K(0;1;0)$ và có véc-tơ pháp tuyến $\overrightarrow{AK}=(-2;0;-1)$ là
			$$-2(x-0)+0(y-1)-1(z-0)=0 \Leftrightarrow 2x+z=0.$$
			Vậy $a+b+c=3$.	
		}{
			\begin{tikzpicture}[line join=round, line cap=round,>=stealth,thick,scale=0.5]
				%Khai báo điểm----------------------
				\path (0,0) coordinate (H);
				\path (0,3) coordinate (A);
				\path (3,-1) coordinate (K);
				\path (-1,1) coordinate (D);	
				\path (-3,-3) coordinate (E);
				\path (5,-3) coordinate (F);	
				\path (5,0) coordinate (I);
				\path (1,-2) coordinate (G);
				\path pic[draw,angle radius=5pt]{right angle= A--K--I};
				\path pic[draw,angle radius=5pt]{right angle= A--H--K};
				\path pic["\scriptsize$P$", angle eccentricity=0.6,draw,angle radius=20pt]{angle= F--E--D};
				%vẽ---------------------------------
				\draw (A)--(H)--(K)--cycle;
				\draw (D)--(E)--(F);
				\draw (I)--(K)--(G);
				\node at (4,0){$Oy$};
				%hiển thị tên điểm----------------------
				\foreach \t/\g in {H/-90,K/-90,A/90}{
					\draw[fill=black] (\t) circle (1pt) node[shift={(\g:10pt)}]{$ \t $};
				}
			\end{tikzpicture}	
		}
	}
\end{ex}


%%%==============EX_7============%%%
\begin{ex}%[2H5V1-5]
	Trong KG $Oxyz$, cho điểm $A(1;2;2)$. Gọi $(P)$ là mặt phẳng chứa trục $Ox$ sao cho khoảng cách từ $A$ đến $(P)$ lớn nhất. Phương trình của mặt phẳng $(P)$ có dạng $y+az=0$. Tính $a$.
	\shortans[0]{$1$}
	\loigiai{
		Gọi hình chiếu vuông góc của điểm $A(1;2;2)$ lên trục $Ox$ là $M(1;0;0)$. \\
		Khoảng cách từ $A$ đến $(P)$ lớn nhất nên mặt phẳng $(P)$ có véctơ pháp tuyến là $\overrightarrow{MA}=(0;2;2)$. \\
		Phương trình mặt phẳng $(P)$ đi qua điểm $M(1;0;0)$ và có véctơ pháp tuyến là $\overrightarrow{MA}=(0;2;2)$ có dạng $0.(x-1)+2(y-0)+2(z-0)=0$ $\Leftrightarrow y+z=0$.
	}
\end{ex}
%%%==============EX_8============%%%
\begin{ex}%[2H5V2-7]
	Trong KG $Oxyz$, cho điểm $A(2;-1;-2)$ và đường thẳng $(d)$ có phương trình $\dfrac{x-1}{1}=\dfrac{y-1}{-1}=\dfrac{z-1}{1}$. Gọi $(P)$ là mặt phẳng đi qua điểm $A$, song song với đường thẳng $(d)$ và khoảng cách từ $(d)$ tới mặt phẳng $(P)$ là lớn nhất. Khi đó mặt phẳng $(P)$ có dạng $x+ay+bz+c=0$. Tính $T=a+b+c+d$.
	\shortans[0]{$-15$}
	\loigiai{
		Gọi $H$ là hình chiếu của $A$ lên đường thẳng $d$. Ta suy ra $H(1;1;1)$. \\
		Gọi $(P)$ là mặt phẳng đi qua điểm $A$ và song song với đường thẳng $(d)$. \\
		Gọi $K$ là hình chiếu của $H$ lên mặt phẳng $(P)$. \\
		Do $d\parallel (P)$ nên ta có $\mathrm{d}\left(d,(P)\right)=\mathrm{d}\left(H,(P)\right)=HK \le HA$. \\
		Như vậy khoảng cách từ đến lớn nhất bằng $HA$. Và khi đó $(P)$ nhận $AH=(-1;2;3)$ làm vectơ pháp tuyến. \\
		Do $(P)$ đi qua $A(2;-1;-2)$ nên ta có phương trình của $(P)$ là $x-2y-3z-10=$ nên $T=a+b+c=-15$.
	}
\end{ex}
%%%%==============EX_9============%%%
\begin{ex}%[2H5C1-5]
	Trong KG $Oxyz$, cho điểm $A(2;5;3)$ và đường thẳng $d\colon \dfrac{x-1}{2}=\dfrac{y}{1}=\dfrac{z-2}{2}$. Gọi $(P)$ là mặt phẳng chứa $(d)$ sao cho khoảng cách từ $A$ đến $(P)$ là lớn nhất. Biết khoảng cách từ gốc tọa độ $O$ đến mặt phẳng $(P)$ bằng $\dfrac{m\sqrt{n}}{n}$ với $m,n\in\mathbb{Z}$. Tính $T=m+n$.
	\shortans[0]{$3$}
	\loigiai{
	Gọi $\overrightarrow{n}=(a,b,c)$ là một vectơ pháp tuyến của $(P)$ với $a^2+b^2+c^2\neq 0$. \\
		Điểm $M(1;0;2)\in d \Rightarrow M\in(P)$. \\
		Phương trình của $(P)\colon ax+by+cz-(a+2c)=0$. \\
		Một vectơ chỉ phương của $(d)$ là\\ $\overrightarrow{u}=(2;1;2)\\ \Rightarrow \overrightarrow{n}\perp \overrightarrow{u} \Leftrightarrow \vec{n}\cdot \vec{u}=0 \Leftrightarrow 2a+b+2c=0 \Leftrightarrow b=-2a-2c$.\\
		Khi đó
		$\mathrm{d}\left(A,(P)\right)=\dfrac{|a+5b+c|}{\sqrt{a^2+b^2+c^2}}=\dfrac{9|a+c|}{\sqrt{a^2+c^2+4(a+c)^2}}$. \\
		Ta có $\dfrac{(a+c)^2}{2} \le a^2+c^2$.\\
		Suy ra $a^2+c^2+4(a+c)^2 \ge \dfrac{(a+c)^2}{2} +4(a+c)^2=\dfrac{9}{2}(a+c)^2$. \\
		Do đó $\mathrm{d}\left(A,(P)\right)=\dfrac{9|a+c|}{\sqrt{a^2+c^2+4(a+c)^2}}\le \dfrac{9|a+c|}{\sqrt{\dfrac{9}{2}(a+c)^2}}=3\sqrt{2}$. \\
		Dấu \lq\lq =\rq\rq\, xảy ra khi $\heva{&a=c \\ &b=-4a}$. Chọn $a=c=1$, suy ra $b=-4$.\\
		Phương trình $(P)\colon x-4y+z-3=0 \Rightarrow \mathrm{d}\left(O,(P)\right)=\dfrac{\sqrt{2}}{2}$. \\
		Vậy $T=1+2=3$.
	}
\end{ex}
%%%%==============EX_10============%%%
\begin{ex}%[2H5V1-5]
	Trong KG $Oxyz$, cho điểm $A(1; 1; 2)$ và mặt phẳng $(P)\colon (m-1)x+y+mz-1=0$, với $m$ là tham số. Biết khoảng cách từ điểm $A$ đến mặt phẳng $(P)$ lớn nhất. Tính $m$.
	\shortans[0]{$5$}
	\loigiai{
		Ta có $\mathrm{d}\left(A;(P)\right)=\dfrac{\left|m-1+1+2m-1\right|}{\sqrt{(m-1)^2+1+m^2}}=\sqrt{\dfrac{(3m-1)^2}{2\bigl(m^2-m+1\bigr)}}$. \\
		Xét $f(m)=\dfrac{(3m-1)^2}{2\bigl(m^2-m+1\bigr)} \Rightarrow f'(m)=\dfrac{(5-m)(3m-1)}{2\bigl(m^2-m+1\bigr)^2}=0\Leftrightarrow \left[\begin{array}{l} {m=\dfrac{1}{3}} \\ {m=5.} \end{array}\right.$ \\
		\begin{center}
			\begin{tikzpicture}
				\tkzTabInit[nocadre, lgt=1, espcl=2]{$m$ /1,$f'$ /1,$f$ /2}{$-\infty$,$\dfrac{1}{3}$,$5$,$+\infty$}
				\tkzTabLine{,-,0,+,0,-,}
				\tkzTabVar{+/$\dfrac{9}{2}$ ,-/ $0$, +/ $\dfrac{14}{3}$ /, -/$\dfrac{9}{2}$ /}
			\end{tikzpicture}
		\end{center}
		Vậy $\max \mathrm{d}\left(A;(P)\right)=\sqrt{\dfrac{14}{3}}$ khi $m=5$.
	}
\end{ex}
%%%%==============EX_11============%%%
\begin{ex}%[2H5V1-5]
	Trong KG $Oxyz$, cho hai điểm $A(1;2;-1),B(3;0;3)$. Gọi $(P)$ là mặt phẳng đi qua điểm $A$ và cách $B$ một khoảng lớn nhất. Biết $\overrightarrow{n}=(1;a;b)$ là một véc-tơ pháp tuyến của $(P)$. Tính $T=a+b$.
	\shortans[0]{$1$}
	\loigiai{
		Ta có $\overrightarrow{AB}=(2;-2; 4)\Rightarrow AB=2\sqrt{6}$. \\
		Gọi $H$ là hình chiếu vuông góc của $B$ trên mặt phẳng $(P)$. \\
		Ta có $\mathrm{d}\left(B,(P)\right)=BH\le BA=2\sqrt{6} \Rightarrow \max \mathrm{d} \left(B,(P)\right)=2\sqrt{6}$, đạt được khi $H\equiv A$. \\
		Khi đó mặt phẳng $(P)$ đi qua $A$ và nhận $\overrightarrow{AB}=(2;-2; 4)$ làm véc-tơ pháp tuyến.
		Suy ra $\overrightarrow{n}=(1;-1;2)$ là một véc-tơ pháp tuyến của mặt phẳng $(P)$. \\
		Vậy $T=1$.
	}
\end{ex}
%%%%==============EX_12============%%%
\begin{ex}%[2H5V1-5]
	Trong KG $Oxyz$, cho điểm $A(1;2;-3)$ và mặt phẳng $(P)\colon 2x+2y-z+9=0$. Đường thẳng $d$ qua và vuông góc với $(Q)\colon 3x+4y-4z+5=0$, cắt mặt phẳng $(P)$ tại $B$. Điểm $M$ nằm trong $(P)$ sao cho $M$ luôn nhìn $AB$ dưới góc vuông. Tính độ dài lớn nhất của $MB$ (làm tròn một chữ số).
	\shortans[0]{$2{,}2$}
	\loigiai{
		Đường thẳng $(d)$ qua $A$ và vuông góc với $(Q)$ có dạng $\heva{&x=1+3t \\ &y=2+4t \\ &z=-3-4t.}$ \\
		Gọi $B$ là giao điểm của $(d)$ và mặt phẳng $(P)$ suy ra $B(-2;-2;1)$. \\
		Do $\widehat{AMB}=90^{\circ}$ nên $M$ thuộc mặt cầu $(S)$ có đường kính $AB$. \\
		Suy ra $M$ thuộc đường tròn $(C)$ là giao tuyến của mặt cầu $(S)$ và mặt phẳng $(P)$. \\
		Suy ra $BM$ lớn nhất khi $BM$ bằng đường kính của đường tròn $(C)$. \\
		Ta có $(S)$ có tâm $I\left(-\dfrac{1}{2};0;-1\right)$, bán kính $R=\dfrac{\sqrt{41}}{2}$. \\
		Ta có $\mathrm{d}\left(I,(P)\right)=3$. \\
		Bán kính của đường tròn $(C)$ là $r=\sqrt{R^2-\mathrm{d}^2\left(I,(P)\right)}=\dfrac{\sqrt{5}}{2}$. \\
		Vậy $BM$ lớn nhất bằng $\sqrt{5}\approx 2{,}2$.
	}
\end{ex}





%%%==============Cau_EX13==============%%%
\begin{ex}%[2H2V2-2]
Cho $A(4;5;6)$; $B(1;1;2)$, $M$ là một điểm di động trên mặt phẳng $(P)\colon 2x+y+2z+1=0$. Tính giá trị lớn nhất của biểu thức $|MA-MB|$ (làm tròn kết quả đến hàng phần trăm).
\shortans[]{$6{,}40$}
\loigiai{
Ta có $|MA-MB|\leq AB$ với mọi điểm $M\in(P)$.\\
Vì $(2\cdot 4+5+2\cdot 6+1) \cdot(2\cdot 1+1+2\cdot 2+1)=208> 0$ nên hai điểm $A, B$ nằm cùng phía với $(P)$.\\
Dấu \lq\lq  =\rq\rq\, xảy ra khi và chỉ khi $M=AB\cap(P)$.\\
Khi đó, $|MA-MB|$ nhận giá trị lớn nhất là $$AB=\sqrt{(4-1)^2+(5-1)^2+(6-2)^2}=\sqrt{41}\approx 6{,}40.$$}
\end{ex}
%%%==============HetCau_EX13==============%%%

%%%==============Cau_EX14==============%%%
\begin{ex}%[2H2V2-2]
Trong KG $Oxyz$, cho ba điểm $A(-1;0;1)$, $B(3;2;1)$, $C(5;3;7)$. Gọi $M(a;b;c)$ là điểm thỏa mãn $MA=MB$ và $MB+MC$ đạt giá trị nhỏ nhất. Tính $P=a+b+c$.
\shortans[]{$5$}
\loigiai{
Gọi $I$ là trung điểm của $AB$, suy ra $I(1;1;1)$ và  $\overrightarrow{AB}=(4;2;0)$.\\
Điểm $M$ thỏa mãn $MA=MB$ khi $M\in(\alpha)$.\\
Phương trình mặt phẳng trung trực của $AB$ là $(\alpha)\colon 2x+y-3=0$.\\
Vì $(2\cdot 3+1\cdot 2-3) \cdot(2\cdot 5+1\cdot 3-3)=50> 0$ nên $B, C$ nằm về một phía so với $(\alpha)$, suy ra $A, C$ nằm về hai phía so với $(\alpha)$.\\
Khi đó $MB+MC=MA+MC\geq AC$.\\
$MB+MC$ nhỏ nhất bằng $AC$ khi $M=AC\cap(\alpha)$.\\
PTĐT $AC\colon\heva{&x=-1+2t \\
	&y=t \\
	&z=1+2t}$, do đó tọa độ điểm $M$ là nghiệm của hệ phương trình $2\cdot (-1+2t)+t-3=0\Leftrightarrow t=1$.\\
Suy ra, tọa độ $M(1;1;3)$. Nên $a+b+c=5$.
}
\end{ex}
%%%==============HetCau_EX14==============%%%

%%%==============Cau_EX15==============%%%
\begin{ex}%[2H2V2-2]
Trong KG $Oxyz$, cho $A(-3;1;1)$, $B(5;1;1)$ và hai mặt phẳng $(P)\colon x+2y+z-4=0$, $(Q)\colon -x+y+z-1=0$. Gọi $M(a;b;c)$ là điểm nằm trên hai mặt phẳng $(P)$ và $(Q)$ sao cho $MA+MB$ đạt giá trị nhỏ nhất. Tính $T=a^2+b^2+c^2$.
\shortans[]{$3$}
\loigiai{
Từ giả thiết ta có $M$ thuộc giao tuyến $d$ của hai mặt phẳng $(P)$ và $(Q)$.
\begin{itemize}
	\item Mặt phẳng $(P)$ có một véctơ pháp tuyến là $\overrightarrow{n}_P=(1;2;1)$.
	\item Mặt phẳng $(Q)$ có một véctơ pháp tuyến là $\overrightarrow{n}_Q=(-1;1;1)$.
\end{itemize}
Khi đó đường thẳng $d$ đi qua $N(1;1;1)$ và có một véctơ chỉ phương là $\overrightarrow{u}_d=\left[\overrightarrow{n}_P, \overrightarrow{n}_Q\right]=(1;-2;3)$ nên có PTTS là $d\colon \heva{&x=1+t \\
	&y=1-2t \\
	&z=1+3t}$ suy ra $M(1+t;1-2t;1+3t)$.\\
$MA+MB=\sqrt{(t+4)^2+4t^2+9t^2}+\sqrt{(t-4)^2+4t^2+9t^2}=\sqrt{14t^2+8t+16}+\sqrt{14t^2-8t+16}$.\\
Bài trở thành tìm giá trị nhỏ nhất của hàm số
$$
\begin{aligned}
		& f(t)=\sqrt{14 t^2+8 t+16}+\sqrt{14 t^2-8 t+16}=\sqrt{14}\left[\sqrt{t^2+\dfrac{4}{7} t+\dfrac{8}{7}}+\sqrt{t^2-\dfrac{4}{7} t+\dfrac{8}{7}}\right] \\
		&=\sqrt{14}\left[\sqrt{\left(t+\dfrac{2}{7}\right)^2+\left(\dfrac{2}{\sqrt{7}}\right)^2}+\sqrt{\left(t-\dfrac{2}{7}\right)^2+\left(\dfrac{2}{\sqrt{7}}\right)^2}\right].
\end{aligned}
$$
Đặt $\overrightarrow{u}=\left(t+\dfrac{2}{7}; \dfrac{2}{\sqrt{7}}\right)$ và $\overrightarrow{v}=\left(\dfrac{2}{7}-t; \dfrac{2}{\sqrt{7}}\right)$.\\
Khi đó\\ $f(t)=\sqrt{14}[|\overrightarrow{u}|+|\overrightarrow{v}|] \geq \sqrt{14}|\overrightarrow{u}+\overrightarrow{v}|$.\\ Suy ra $f(t) \geq \sqrt{14} \cdot \sqrt{\left(\dfrac{4}{7}\right)^2+\left(\dfrac{4}{\sqrt{7}}\right)^2}=\dfrac{8\sqrt{7}}{49}$.\\
Dấu \lq\lq =\rq\rq\, xảy ra khi và chỉ khi hai véctơ $\overrightarrow{u}$ và $\overrightarrow{v}$ cùng hướng hay $\dfrac{t+\dfrac{2}{7}}{\dfrac{2}{7}-t}=\dfrac{\dfrac{2}{\sqrt{7}}}{\dfrac{2}{\sqrt{7}}} > 0\Rightarrow t=0$.\\
Do đó $M(1;1;1)$.\\
Vậy $T=a^2+b^2+c^2=3$.
	}
\end{ex}
%%%==============HetCau_EX15==============%%%

%%%==============Cau_EX16==============%%%
\begin{ex}%[2H2V2-2]
Trong KG $Oxyz$ cho mặt phẳng $(P)\colon y-1=0$, đường thẳng $d\colon \left\{\begin{array}{l}x=1\\ y=2-t \\ z=1\end{array}\right.$ và hai điểm $A(-1;-3; 11), B\left(\dfrac{1}{2}; 0; 8\right)$. Hai điểm $M, N$ thuộc mặt phẳng $(P)$ sao cho $d(M, d)=2$ và $NA=2NB$. Tìm giá trị nhỏ nhất của đoạn $MN$.
\shortans[]{$1$}
\loigiai{
Gọi $I=d \cap(P) \Rightarrow I(1;2-t;1)$.\\
Ta có $I\in(P) \Rightarrow 2-t-1=0\Rightarrow t=1\Rightarrow I(1;1;1)$.\\
Ta lại có $d \perp(P) \Rightarrow M$ thuộc đường tròn tâm $I(1;1;1)$, $R_1=2$.\\
Gọi $N(x;y;z) \Rightarrow \heva{&\overrightarrow{N A}(-1-x;-3-y;11-z)\\& \overrightarrow{N B}\left(\dfrac{1}{2}-x;-y;8-z\right).}$\\
Mà \begin{eqnarray*}
& & NA=2NB\\
&\Leftrightarrow&(1+x)^2+(3+y)^2+(11-z)^2=4\left[\left(\dfrac{1}{2}-x\right)^2+y^2+(8-z)^2\right]\\
&\Leftrightarrow& 3 x^2+3 y^2+3 z^2-6 x-6 y-42 z+126=0\\
&\Leftrightarrow& x^2+y^2+z^2-2 x-2 y-14 z+42=0.
\end{eqnarray*} 
Vậy $N\in S\left(J(1;1;7);R_2=3\right)$ và $J\in(P)\colon y=1$.\\
Nên $N$ thuộc đường tròn tâm $J(1;1;7)$; $R_2=3$.\\
Ta có $IJ=6> R_1+R_2\Rightarrow MN_{\min}=IJ-R_1-R_2=1$.
}
\end{ex}
%%%==============HetCau_EX16==============%%%

%%%==============Cau_EX17==============%%%
\begin{ex}%[2H2V2-2]
Trong KG $Oxyz$, cho bốn điểm $A(2;0;1)$, $B(3;1;5)$, $C(1; 2; 0)$, $D(4; 2; 1)$. Gọi $(\alpha)$ là mặt phẳng đi qua $D$ sao cho ba điểm $A$, $B$, $C$ nằm cùng phía đối với $(\alpha)$ và tổng khoảng cách từ các điểm $A$, $B$, $C$ đến mặt phẳng $(\alpha)$ là lớn nhất. Giả sử phương trình $(\alpha)$ có dạng $2x+m y+n z-p=0$. Khi đó, tính giá trị của $T=m+n+p$.
\shortans[]{$9$}
\loigiai{
Vì mặt phẳng $(\alpha)$ đi qua $D(4;2;1)$ nên phương trình $(\alpha)$ có dạng
	$a \cdot(x-4)+b \cdot(y-2)+c \cdot(z-1)=0$.\\
Đặt $S=\mathrm{d}[A,(\alpha)]+\mathrm{d}[B,(\alpha)]+\mathrm{d}[C,(\alpha)]=\dfrac{|-2a-2b|+|-a-b+4c|+|-3a-c|}{\sqrt{a^2+b^2+c^2}}$.\\
Theo giả thiết, $A$, $B$, $C$ nằm cùng phía đối với $(\alpha)$ nên không mất tính tổng quát, ta giả sử
	$$
\heva{&-2a-2b > 0\\
	&-a-b+4c > 0\\
	&-3a-c > 0.}
	$$
Khi đó, $S=\dfrac{-2a-2b-a-b+4c-3a-c}{\sqrt{a^2+b^2+c^2}}=\dfrac{-6a-3b+3c}{\sqrt{a^2+b^2+c^2}}$.\\
Áp dụng bất đẳng thức Cauchy-Schwars cho hai bộ số $(-6;-3; 3)$ và $(a; b; c)$, ta được
	$$-6a-3b+3c \leq|-6a-3b+3c|\leq \sqrt{\left(6^2+3^2+3^2\right) \cdot\left(a^2+b^2+c^2\right)} \Rightarrow S\leq 3\sqrt{6}.$$
	Đẳng thức xảy ra $\Leftrightarrow\heva{&-6a-3b+3c \geq 0\\	&\dfrac{a}{-6}=\dfrac{b}{-3}=\dfrac{c}{3}.}$\\
	Ta chọn$\heva{&a=-2\\
		&b=-1\\
		&c=1.}$\\
	$\Rightarrow(\alpha)\colon-2x-y+z+9=0$ hay $(\alpha)\colon 2x+y-z-9=0\Rightarrow m=1, n=-1, p=9$.\\
	Vậy $T=m+n+p=9$.
	}
\end{ex}
%%%==============HetCau_EX17==============%%%